\section{Related Work}
\label{sec:related_work}

\para{Classical weight decay.} The classical argument for L2 regularization is that smaller-norm solutions can improve generalization by reducing variance and limiting sensitivity to target noise \cite{krogh1991simple}. That view motivates our core expectation that L2 should shrink coefficients and reduce the train-validation gap on small datasets. Our study directly tests that practical implication in a simple linear model rather than in a high-capacity deep network.

\para{Optimizer semantics and implementation details.} Later work showed that L2 regularization and weight decay are not always interchangeable. Loshchilov and Hutter \cite{loshchilov2019decoupled} demonstrate that the equivalence breaks under adaptive optimizers, which matters whenever a comparison loosely uses ``weight decay'' to mean several different update rules. Our setting avoids that ambiguity by using the same \texttt{lbfgs} logistic-regression solver in both conditions and changing only the penalty strength.

\para{Regularization is not the whole story.} Zhang et al.\ \cite{zhang2017understanding} caution against explanations of generalization that rely too heavily on explicit regularizers. That caution is directly relevant here: even if L2 reduces overfitting indicators, it need not deliver a reliable test-accuracy gain. We therefore treat predictive improvement and gap reduction as separate empirical questions instead of assuming they move together.

\para{Small-data and dynamics-oriented views.} Recent work strengthens the case for studying regularization in low-data regimes while still keeping the claim conditional. Power et al.\ \cite{power2022grokking} show that small datasets can exhibit long delays between memorization and generalization, with regularization affecting data efficiency. Golatkar et al.\ \cite{golatkar2019time} argue that the timing of regularization matters, especially early in training. D'Angelo et al.\ \cite{dangelo2024weight} further reinterpret weight decay as a mechanism that shapes optimization dynamics rather than only limiting capacity. These papers motivate our small-data focus, but they also suggest that predictive gains may depend strongly on dataset structure and training conditions.

\para{Position of this work.} Unlike prior deep-learning studies, we focus on a narrow but common practical setting: tiny real tabular classification datasets paired with logistic regression. Unlike purely theoretical treatments, we quantify both the size and the consistency of the observed L2 effect across repeated stratified splits. Our contribution is therefore empirical and measurement-oriented: we document when shrinkage reliably appears and when predictive gains fail to rise above split noise.
